\section*{الفصل \ref{ch:g12:reaction-rates} --- سرعات التفاعلات}

\begin{solution}{exo:g12:reaction-rates:1}
درجة الحرارة؛ مساحة التماس؛ التركيز؛ درجة الحرارة.
\end{solution}

\begin{solution}{exo:g12:reaction-rates:2}
\qty{13.9}{min} (بالأزرق) و\qty{6.9}{min} (بالأحمر). وبعد زمني نصف تفاعل،
$10.0 / 4 = \qty{2.5}{mmol/L}$.
\end{solution}

\begin{solution}{exo:g12:reaction-rates:3}
الميل $(5.0 - 1.0)/10 = \qty{0.40}{mmol.L^{-1}.min^{-1}}$: وهذه هي
\omterm{def:g12:reaction-rates:rate}{سرعة الظهور} عند \qty{5}{min}.
\end{solution}

\begin{solution}{exo:g12:reaction-rates:4}
التبريد \omterm{def:g10:concentration-and-dilution:dilution}{والتخفيف} يجعلان التفاعل بطيئًا إلى حد أن التركيب لا
يعود يتغير خلال \omterm{def:g11:titration:titration}{المعايرة}. والتركيب المقيس هو التركيب الذي
كان \omterm{def:g6:pure-substances-and-mixtures:mixture}{للخليط} حين أُوقف التفاعل، أي لحظة الصب.
\end{solution}

\begin{solution}{exo:g12:reaction-rates:5}
\qty{30}{min} تساوي 3 أزمنة نصف تفاعل: فيبقى $1/8$. و\qty{1}{h} تساوي 6 أزمنة نصف تفاعل:
$1/64$، أي نحو \qty{1.6}{\percent}.
\end{solution}

\begin{solution}{exo:g12:reaction-rates:6}
$\ln[\ce{A}]$: 2.079، 1.733، 1.386، 1.040، 0.693. وهو ينخفض بمقدار 0.346 كل
\qty{5}{min}: خط مستقيم، فالرتبة أولى. $k = 0.346/5 =
\qty{0.0693}{min^{-1}}$ و$t_{1/2} = \ln 2 / k = \qty{10.0}{min}$ (كما
يُرى مباشرة: 8.00، 4.00، 2.00 كل \qty{10}{min}).
\end{solution}

\begin{solution}{exo:g12:reaction-rates:7}
$k = 0.693 / 25 = \qty{0.0277}{s^{-1}}$.
\end{solution}

\begin{solution}{exo:g12:reaction-rates:8}
$0.100\,e^{-0.20} = \qty{0.0819}{mol/L}$؛ $0.100\,e^{-0.60} =
\qty{0.0549}{mol/L}$؛ $0.100\,e^{-2.0} = \qty{0.0135}{mol/L}$.
\end{solution}

\begin{solution}{exo:g12:reaction-rates:9}
عند \qty{13.9}{min}، $[\ce{A}] = \qty{5.0}{mmol/L}$؛ وميل المماس
قريب من $-0.25$، و$-k[\ce{A}] = -0.050 \times 5.0 =
\qty{-0.25}{mmol.L^{-1}.min^{-1}}$: أي نصف القيمة الابتدائية.
\end{solution}

\begin{solution}{exo:g12:reaction-rates:10}
ثنائي اليود هو النوع الملوّن الوحيد: فبحسب قانون بير--لامبرت
$A = \varepsilon \ell [\ce{I2}]$. \omterm{def:g12:reaction-rates:rate}{وسرعة الظهور} هي ميل
مماس المنحنى $[\ce{I2}](t)$، أي ميل
مماس $A(t)$ مقسومًا على $\varepsilon \ell$.
\end{solution}

\begin{solution}{exo:g12:reaction-rates:11}
السرعة الابتدائية $k[\ce{A}]_0$ أكبر بمرتين في التجربة الثانية؛
وزمنا نصف التفاعل متساويان، $\ln 2 / k$، ولا يتوقفان على $[\ce{A}]_0$.
\end{solution}

\begin{solution}{exo:g12:reaction-rates:12}
$t = \ln(1/f)/k$. وبأزمنة نصف التفاعل: $\ln(1/f)/\ln 2$. ولنسبة \qty{1}{\percent}،
$\ln 100 / \ln 2 = 6.6$؛ ولنسبة \qty{0.1}{\percent}، $\ln 1000 / \ln 2 =
10.0$.
\end{solution}

\begin{solution}{exo:g12:reaction-rates:13}
استهلاك \qty{90}{\percent}، أي $f = 0.10$: $t = \ln 10 / k$، أي
\qty{46}{min} عند \qty{20}{\celsius} و\qty{23}{min} عند \qty{30}{\celsius}.
\omterm{def:g12:reaction-rates:first-order}{ثابت السرعة} يتضاعف على مدى \qty{10}{\celsius}: فالقاعدة صحيحة هنا.
\end{solution}

\begin{solution}{exo:g12:reaction-rates:14}
تنخفض درجة الحرارة إلى قرب \qty{0}{\celsius}، وتُقسم التراكيز
على 10. ومع سرعة متناسبة مع جداء
تركيزين، يقسمها \omterm{def:g10:concentration-and-dilution:dilution}{التخفيف} وحده على $10 \times 10 = 100$.
\end{solution}

\begin{solution}{exo:g12:reaction-rates:15}
$v(t) = -\dfrac{d}{dt}\big([\ce{A}]_0 e^{-kt}\big) = k[\ce{A}]_0 e^{-kt}$.
السرعة الابتدائية $k[\ce{A}]_0$؛ وتنتصف السرعة حين $e^{-kt} = 1/2$، أي
بعد زمن نصف تفاعل واحد.
\end{solution}

\begin{solution}{pb:g12:reaction-rates:1}
\textbf{1.} بحسب قانون بير--لامبرت، إذ الصبغة هي النوع الوحيد
الذي يمتص عند هذا الطول الموجي؛ ويجب أن يكون \omterm{def:g2:mixing-and-dissolving:dissolve}{المحلول} مخففًا بما يكفي
($A$ أقل من نحو 1).

\textbf{2.} لا يمتص أي ضوء مرئي: فلا يضيف شيئًا إلى $A$.

\textbf{3.} لكي لا يكاد تركيزه يتغير: فتتوقف السرعة عندئذ
على الصبغة وحدها.

\textbf{4.} $A = 0$ (الشاهد).

\textbf{5.} تنخفض $A$ من 0.800 إلى 0.402 في \qty{6}{min}: أي نحو
\qty{6}{min}.

\textbf{6.} من 0.402 عند \qty{6}{min} إلى 0.199 عند \qty{12}{min}: تنتصف
من جديد في \qty{6}{min}.

\textbf{7.} $-0.223$، $-0.448$، $-0.691$، $-0.911$، $-1.155$، $-1.366$،
$-1.614$، $-1.945$، $-2.551$، $-3.079$.

\textbf{8.} تقع النقاط على خط مستقيم: \omterm{def:g12:reaction-rates:first-order}{فالتفاعل من الرتبة الأولى}
بالنسبة إلى الصبغة.

\textbf{9.} $(-2.551 + 0.223)/20 = \qty{-0.116}{min^{-1}}$.

\textbf{10.} $k = \qty{0.116}{min^{-1}}$.

\textbf{11.} $0.693 / 0.116 = \qty{6.0}{min}$، كما قُرئ في الجدول.

\textbf{12.} $k A_0 = 0.116 \times 0.800 = \qty{0.093}{min^{-1}}$ بوحدات
\omterm{def:g11:absorbance:absorbance}{الامتصاصية} في الدقيقة.

\textbf{13.} القيمة نفسها، \qty{6.0}{min}: فزمن نصف \omterm{def:g12:reaction-rates:first-order}{التفاعل من الرتبة الأولى} لا
يتوقف على القيمة الابتدائية.

\textbf{14.} $0.800\, e^{-0.116 \times 30} = 0.025$.

\textbf{15.} $t = \ln 100 / k = 4.61 / 0.116 = \qty{40}{min}$.

\textbf{16.} $40 / 6.0 = 6.6$ من أزمنة نصف التفاعل.

\textbf{17.} $t = \ln(0.800/0.010)/k = 4.38 / 0.116 = \qty{38}{min}$.

\textbf{18.} دائمًا 6.6 من أزمنة نصف التفاعل، أي الآن $6.6 \times 3.0 = \qty{20}{min}$.

\textbf{19.} نحو 6.6 من أزمنة نصف التفاعل، أي نحو \qty{40}{min} هنا.
\end{solution}
